\documentclass[UTF8,11pt,fontset=fandol]{ctexart}
\usepackage[a4paper,margin=25mm]{geometry}
\usepackage{amsmath,booktabs,longtable,array,hyperref,fancyhdr}
\usepackage{fontspec}
\usepackage{xurl}
\usepackage{unicode-math}
\setmainfont{STIXTwoText}[Extension=.otf,
  UprightFont=*-Regular, BoldFont=*-Bold,
  ItalicFont=*-Italic, BoldItalicFont=*-BoldItalic]
\setmathfont{STIXTwoMath-Regular.otf}
\linespread{1.3}
\setlength{\parskip}{4pt}
\setlength{\emergencystretch}{3em}
\hypersetup{colorlinks=true,linkcolor=black,urlcolor=blue}
\pagestyle{fancy}
\fancyhf{}
\fancyhead[L]{\small 刘星辰}
\fancyhead[R]{\small 2026.11}
\fancyfoot[C]{\small\thepage}
\renewcommand{\headrulewidth}{0.4pt}
\begin{document}
\begin{center}
{\LARGE\bfseries 哈密顿路极值序列的非代数性与公式限制\par}
\vspace{1.2em}
{\large 刘星辰\par}
\vspace{0.5em}
{\normalsize 个人研究者\par}
{\normalsize \texttt{lxc-em5158@outlook.com}\par}
\vspace{0.5em}
{\small 2026年11月10日 · 本地修订版 v0.4.0\par}
\end{center}
\vspace{1em}

\section*{摘要}
记 $P(n)$ 为 $n$ 阶竞赛图中有向哈密顿路数的最大值。本文研究其精确表达所受到的限制。利用已有的渐近定理，证明与 $P(n)$ 相关的三个生成函数均非代数。另由有理序列的幂次子序列极限，排除任意固定进制的有限维线性数字表示。在明确的边界条件R下，进一步排除若干有限阶线性递推和伸缩关系。本文不证明实际极值序列满足R，因此后一组结论保留条件性。正文介绍结果的含义与证明思路，完整定义和证明列于附录。这些结果仅限制特定的表达形式，不能据此断言一般闭式不存在，也不构成计算困难性的证明。

关键词：竞赛图；哈密顿路；极值计数；生成函数；闭式限制。



\section{问题与研究思路}
竞赛图中，任意两个不同顶点之间恰有一条有向边。有向哈密顿路沿边的方向恰好经过每个顶点一次。设 $H(T)$ 为竞赛图 $T$ 中构成这种路径的顶点排列数，定义

\[
P(n)=\max_{|V(T)|=n}H(T).
\]
为研究 $P(n)$ 的精确表达，本文首先考察其中几类可以排除的公式。

为分离 $P(n)$ 的主要增长项，取随机竞赛图的平均路径数作为归一化基准：

\[
\mu_n=\frac{n!}{2^{n-1}}.
\]
共有 $n!$ 个顶点排列，每个排列中连续的 $n-1$ 条边均具有所需方向的概率为 $2^{-(n-1)}$，由此得到上述平均值。引入归一化量

\[
a_n=\frac{P(n)}{\mu_n}.
\]
归一化量 $a_n$ 表示最大路径数与随机平均值之比。Divoux与Wang的渐近定理[1, Theorem 1.4]给出

\[
a_n\longrightarrow L=\frac{\cosh1}{\cos1}.
\]
本文以这一已有极限为出发点，不再重证。

定义数列 $(a_n)$ 的生成函数

\[
A(z)=a_1z+a_2z^2+\cdots=\sum_{n\ge1}a_nz^n.
\]
生成函数把整个数列表示为一个函数，便于研究它能满足哪些方程。对生成函数表达形式的限制，并不等于对所有通项公式的排除。

\section{不依赖结构分类的排除结论}
若函数满足某个非零多项式方程 $F(z,G(z))=0$，则称其为代数函数。有理函数以及平方根等代数表达均属于这一类。

\textbf{结论1.} $A(z)$ 不是代数函数。原始极值序列的指数生成函数

\[
E(z)=\sum_{n\ge1}\frac{P(n)}{n!}z^n=2A(z/2)
\]
也不是代数函数。即使允许方程的系数取任意复数，结论仍然成立。

证明利用两个事实：常数 $L=\cosh1/\cos1$ 是超越数，即它不是任何非零有理系数多项式的根；各项 $a_n$ 则均为有理数。若 $A$ 是代数函数，由于其幂级数系数均为有理数，下述极限在存在且有限时必须为代数数：

\[
\lim_{z\to1^-}(1-z)A(z).
\]
由已知渐近结果，该极限等于 $L$，与其超越性矛盾。完整论证见附录A。

下面考察一个将规模 $n$ 与规模 $n/2$ 联系起来的递推形式，并研究其余项。对实际序列 $(a_n)$，在偶数 $n$ 处定义 $m_n$，使

\[
n^2a_n=2a_{n/2}+m_n.
\]
奇数时定义 $m_n=n^2a_n$。余项的生成函数为

\[
M(z)=\sum_{n\ge1}m_nz^n.
\]
\textbf{结论2.} $M(z)$ 也不是代数函数。因此，这种半规模递推的余项不可能具有有理或代数生成函数。

证明沿用结论1的思路，此时使用的极限为

\[
\lim_{z\to1^-}(1-z)^3M(z)=2L.
\]
结论2可用于筛选上述递推形式的候选余项：若其生成函数为有理函数或代数函数，就不可能恢复实际的 $P(n)$。这一限制仅针对余项的表达形式，并未排除该递推本身。

\textbf{结论3.} 对任意固定整数进制 $q\ge2$，序列 $a_n$、$m_n$ 及 $P(n)$ 均不是复数域上的 $q$-regular序列。这一结论不需要R。它排除固定维度、按输入数字用固定矩阵更新、最后线性读出的精确模型，不排除非线性方法或一般算法。

关键不只是数列的增长速度：对有理值的 $q$-regular序列，若沿 $q$ 的幂次取样并作幂次归一化后存在极限，该极限必须有理。实际 $a_n$ 和 $m_n/n^2$ 的极限是超越数 $L$，从而产生矛盾。完整证明见定理A.5及推论A.6；$P(n)$ 则已由阶乘增长排除。

\section{附加条件下的进一步结论}
进一步的限制取决于 $A(z)$ 在单位圆上的边界行为。记所需的附加条件为R：在单位圆上一组越来越密的点处，某一阶导数首次出现指定强度的对数发散。涉及的边界点、导数阶数和发散系数见附录B；该附录还给出条件R成立的一个充分判据。

\textbf{本节的结论都以实际 $A(z)$ 满足R为前提。本文尚未证明这个前提。}

若R成立，单位圆就是 $A$ 的自然边界，即 $A$ 不能从圆盘内部解析延拓穿过圆周上的任何一点。这种边界行为限制了函数可以满足的有限方程。

利用R规定的更精确边界行为，还可得到以下两类排除结果：

- \textbf{固定阶的多项式系数线性递推。} $P(n)$ 不可能最终满足一个非平凡关系 $\sum_{j=0}^{d}p_j(n)P(n+j)=0$，其中阶数 $d$ 固定，各 $p_j$ 都是多项式。这不排除其他种类的递推。

- \textbf{若干伸缩公式。} 这类公式包括由 $A(z)$、$A(z^2)$、$A(z^3)$ 等函数构成的有限线性关系。若关系不含导数、系数为 $z$ 的有理函数，且清除分母后右端能解析延拓穿过单位圆上除有限点之外的每一点，则该关系不可能成立。允许导数时，若整数基数不是2的幂，仍可排除相应的一整类关系；二进伸缩关系则未被全部排除。完整范围见附录C。


这两类条件性结论分别涉及邻近项之间的线性递推与生成函数的整数伸缩，其适用范围限于上述形式。

\section{结论的适用范围}
生成函数非代数，不能推出其系数数列没有简单通项。例如，数列 $1/n$ 具有简单的通项，而其生成函数 $-\log(1-z)$ 并非代数函数。

自然边界也不排除简短的递推。附录E给出一个不同于实际极值序列的例子：奇数下标取零，偶数下标满足

\[
s_n=-\frac1n+\frac{2}{n^2}s_{n/2}.
\]
该数列的生成函数具有自然边界，而 $s_n$ 可以通过反复减半准确求得。因此，自然边界与简短的求值规则并不矛盾。

上述结果的适用范围如下：

- \textbf{已经证明：} 实际的 $A,E,M$ 均不是代数函数；$a_n,m_n,P(n)$ 在任意固定进制下均不是 $q$-regular。这些结论不需要R或额外的极值图分类。

- \textbf{仍以R为前提：} 上一节列出的递推与伸缩关系限制，以及实际生成函数的自然边界。本文没有证明实际序列满足R。

- \textbf{本文仍未解决：} $P(n)$ 是否有其他形式的精确闭式，是否存在高效的精确算法，以及实际余项 $M$ 是否有可独立求值的超越表达。

本文未得到NP困难性或一般算法下界。把条件结论用于实际极值序列，仍需独立证明R或附录B所列的充分条件。

渐近式与整数性本身也不能推出自然边界。令 $Q_n=\lfloor L\mu_n\rfloor$，$b_n=Q_n/\mu_n$。这些 $Q_n$ 是整数，而 $b_n\in\mathbb Q$ 且 $|b_n-L|\le2^{n-1}/n!$。因此

\[
\sum_{n\ge1}b_nz^n=\frac{Lz}{1-z}+H(z),
\]

其中 $H$ 为整函数。这个生成函数可延拓穿过除 $1$ 外的每个单位圆点，没有单位圆自然边界。它不是实际的竞赛图极值序列，但说明仅凭已有渐近式和同类整数性不足以证明R。


\section*{附录A. 无条件结论的完整证明}
本附录证明正文的结论1—3。沿用正文的 $A,E,M$，并记 $\Theta=z\frac{d}{dz}$。

\textbf{引理A.1.} $L=\cosh1/\cos1$ 为超越数。

\textbf{证明.} 若 $L$ 代数，则

\[
e+e^{-1}-Le^i-Le^{-i}=0.
\]
指数 $1,-1,i,-i$ 是四个不同代数数，系数 $1,1,-L,-L$ 为不全为零的代数数。这违反Lindemann—Weierstrass定理[2]。由于 $\cos1\ne0$，上述表达有定义。证毕。

\textbf{引理A.2（代数关系降域）.} 若 $B(z)\in\mathbb Q[[z]]$ 在 $\mathbb C(z)$ 上代数，则它在 $\mathbb Q(z)$ 上代数。

\textbf{证明.} 清分母后，取非零关系 $H(z,B(z))=0$，其中 $H\in\mathbb C[z,Y]$。对 $H$ 的有限个复系数取一个有理线性基 $\beta_1,\ldots,\beta_s$，写成

\[
H(z,Y)=\sum_{j=1}^s\beta_jH_j(z,Y),\qquad H_j\in\mathbb Q[z,Y].
\]
每个 $H_j(z,B(z))$ 的幂级数系数均为有理数。逐项比较并使用 $\beta_j$ 的有理线性独立性，得到所有 $H_j(z,B(z))=0$。至少一个 $H_j$ 非零，给出所需关系。证毕。

\textbf{引理A.3（有限径向极限）.} 设 $B\in\mathbb Q[[z]]$ 在单位圆盘解析，且在 $\mathbb C(z)$ 上代数。若某个整数 $d\ge0$ 使 $(1-\rho)^dB(\rho)$ 在 $\rho\uparrow1$ 时有有限极限 $c$，则 $c$ 代数。

\textbf{证明.} 由引理A.2，$C(z)=(1-z)^dB(z)$ 满足非零关系 $R(z,C(z))=0$，$R\in\mathbb Q[z,Y]$。除去 $R$ 作为 $Y$ 的多项式时各系数共有的最高次 $(z-1)$ 因子，可假定 $R(1,Y)\not\equiv0$。径向取极限得 $R(1,c)=0$，故结论成立；若 $R(1,Y)$ 是非零常数，则已直接矛盾。证毕。

\textbf{定理A.4.} $A(z)$、$E(z)$ 以及

\[
M(z)=\Theta^2A(z)-2A(z^2)
\]
在 $\mathbb C(z)$ 上均为超越函数。

\textbf{证明.} $a_n$ 有理且趋于 $L$，所以Abel极限给 $(1-\rho)A(\rho)\to L$。引理A.1与A.3排除 $A$ 代数；变量缩放关系随即排除 $E$ 代数。

写 $M(z)=\sum_{n\ge1}m_nz^n$，则

\[
m_n=n^2a_n-2\mathbf1_{\{2\mid n\}}a_{n/2},\qquad m_n\in\mathbb Q,\qquad \frac{m_n}{n^2}\longrightarrow L.
\]
这里指示项为零时无需定义半整数下标。利用

\[
\sum_{n\ge1}n^2\rho^n=\frac{\rho(1+\rho)}{(1-\rho)^3}
\]
得到 $(1-\rho)^3M(\rho)\to2L$。具体而言，将 $m_n/n^2-L$ 分成有限前缀和绝对值小于任意给定 $\varepsilon$ 的尾部，有限前缀趋于零，尾部极限上界为 $2\varepsilon$。再令 $\varepsilon\downarrow0$ 即可。引理A.3再次与 $L$ 超越矛盾。证毕。

因此具体方程 $\Theta^2A-2A(z^2)=F$ 不可能有代数右端，即使代数关系允许任意复常数。非零有理函数倍乘或加上代数函数也不能使 $M$ 代数。但这不排除 $M$ 为超越的D-finite函数。例如 $e^z/(1-z)^3$ 具有有理Taylor系数、是D-finite函数，其归一化径向极限却是超越数 $e$。

\textbf{定理A.5（有理序列的数字表示与极限）.} 设 $q\ge2$ 为整数，$d\ge0$ 为整数，且 $u_n\in\mathbb Q$。若 $(u_n)$ 在 $\mathbb C$ 上为 $q$-regular，且 $q^{-de}u_{q^e}\to c$ 为有限极限，则 $c\in\mathbb Q$。特别地，若 $u_n/n^d\to c$ 且 $c$ 不为有理数，则 $(u_n)$ 不是 $q$-regular。

\textbf{证明.} $q$-kernel的有限维张成空间在取子序列 $n\mapsto qn+j$ 下不变。选取该空间的一组基，将基序列在 $n$ 处的值组成列向量 $v(n)$，便有固定复矩阵 $T_j$ 和行向量 $\ell$，满足 $v(qn+j)=T_jv(n)$ 及 $u_n=\ell v(n)$。因此

\[
u_{q^e}=\ell T_0^e v(1),\qquad B(t):=\sum_{e\ge0}q^{-de}u_{q^e}t^e=\ell(I-tT_0/q^d)^{-1}v(1)\in\mathbb C(t).
\]

虽然表示矩阵可以含任意复常数，$B$ 的Taylor系数均有理，故实际上 $B\in\mathbb Q(t)$。为证明这一点，写 $QB=P$，其中 $P,Q\in\mathbb C[t]$ 且 $Q\ne0$。取所有多项式系数所张成的有限维有理向量空间的一组基 $\beta_i$，写 $P=\sum_i\beta_iP_i$、$Q=\sum_i\beta_iQ_i$，其中 $P_i,Q_i\in\mathbb Q[t]$。逐项比较形式幂级数系数，得到 $Q_iB=P_i$；选取一个非零的 $Q_i$ 即得结论。

由于 $B$ 的系数趋于 $c$，它在单位圆盘内解析，Abel极限定理给出 $(1-t)B(t)\to c$（$t\uparrow1$）。有理系数的有理函数 $(1-t)B(t)$ 在有理点 $t=1$ 具有有限极限，约去可去因子后，该极限必为有理数。于是 $c\in\mathbb Q$。若只需排除超越的 $c$，也可直接将引理A.3以指数1应用于 $B$。证毕。

\textbf{推论A.6.} 对每个整数 $q\ge2$，实际序列 $(a_n)$、$(m_n)$ 和 $(P(n))$ 均不是在 $\mathbb C$ 上的 $q$-regular序列；任意指定有理的第零项不影响这些结论。

\textbf{证明.} 对 $a_n$ 取 $d=0$，对 $m_n$ 取 $d=2$。由 $a_n\to L$ 和 $m_n/n^2\to L$，定理A.5与引理A.1立即排除两种表示。这不使用R。对 $P(n)$，任意固定维度的数字矩阵表示在处理 $n$ 时只作 $O(\log n)$ 次矩阵乘法。选择固定矩阵范数，便得输出的绝对值至多为 $C n^K$，其中 $C,K$ 与 $n$ 无关。另一方面，对随机竞赛图中的哈密顿路数取期望可得 $P(n)\ge\mu_n=n!/2^{n-1}$，其增长快于每个固定幂，矛盾。证毕。

\section*{附录B. 径向条件R及其充分条件}
记 $\nu_2(n)$ 为正整数 $n$ 中因子2的次数，例如 $\nu_2(12)=2$。

以下令 $G$ 为单位圆盘内解析的函数，$C\ne0$ 为固定复数。

\textbf{假设R.} 存在整数 $r_0\ge1$，使对每个固定 $r\ge r_0$ 和每个原始 $2^{r+1}$ 次单位根 $\omega$，函数 $x\mapsto G(\omega x)$ 的所有低于 $2r$ 阶的导数在 $x\uparrow1$ 时有有限极限，并且

\[
\frac{d^{2r}}{dx^{2r}}G(\omega x)=C\,2^{r^2-1}\log(1-x)+O_{r,\omega}(1).
\]
这里是对实际导数的假设，不能通过对未知的大 $O$ 余项形式求导来代替。每次先固定 $r,\omega$，再取径向极限；不要求增长深度的统一界。

\textbf{命题B.1（周期渐近判据）.} 设 $G(z)=\sum_{n\ge1}b_nz^n$，$b_n$ 有界。以下两项足以推出假设R。

- 对每个固定 $s\ge0$ 和任意固定 $J$，在 $\nu_2(n)=s$ 的整数上存在常系数展开 $b_n=\sum_{j=0}^J\alpha_{s,j}n^{-j}+O_{s,J}(n^{-J-1})$。

- 对每个固定 $r\ge r_0$，存在在零点解析的 $B_r$，使对所有充分大的 $2^r\mid n$ 统一有

\[
b_n=B_r(1/n)-\kappa_r\mathbf1_{\{2^{r+1}\mid n\}}n^{-2r-1}+O_r(n^{-2r-2}),\qquad \kappa_r=C\,2^{r(r+1)}.
\]
\textbf{证明.} 固定 $r$，令 $q=2^{r+1}$，$d=2r+1$。在有限个赋值类 $s<r$ 使用第一项，在其余整数使用第二项，拼成

\[
b_n=\sum_{j=0}^{d}c_j(n)n^{-j}+O_r(n^{-d-1}),
\]
其中每个 $c_j$ 为 $q$ 周期。对 $j<d$，它的周期实际整除 $2^r$；最高项可写成 $c_d(n)=u(n)-\kappa_r\mathbf1_{\{q\mid n\}}$，其中 $u$ 也是 $2^r$ 周期。因而对任意原始 $q$ 次根 $\omega$，采用

\[
\widehat c_j(\omega)=\frac1q\sum_{t=1}^{q}c_j(t)\omega^t
\]
的约定，有 $\widehat c_j(\omega)=0$（$j<d$），而 $\widehat c_d(\omega)=-\kappa_r/q$。

有限Fourier分解把显式部分写成多重对数 $\operatorname{Li}_j(z/\zeta)$ 的有限线性组合；非正指标可按几何级数的导数解释。除 $\zeta=\omega$ 外，这些函数在 $z=\omega$ 附近均有与盘内分支相容的解析延拓。唯一首共振项为 $-(\kappa_r/q)\operatorname{Li}_{2r+1}(z/\omega)$。

余项系数为 $O_r(n^{-2r-2})$，求 $2r$ 阶导数后绝对可和，所以该导数有有限径向极限。另一方面，$\operatorname{Li}_{2r+1}(x)$ 的第 $2r$ 阶导数为 $-\log(1-x)+O(1)$，低阶导数均有限。这可由 $\Theta^{2r}\operatorname{Li}_{2r+1}=\operatorname{Li}_1=-\log(1-x)$ 及Euler算子与普通导数的三角关系直接得到。故首系数为 $\kappa_r/q=C2^{r^2-1}$，即R。有限初始项只增加多项式。证毕。

\textbf{与研究笔记的关系.} 固定在提交 \texttt{c049e040} 的研究笔记[7]确实给出了通向实际序列R的证明路线：\texttt{round96core} 中的共同展开与共振跳跃命题F/J对 $r\ge4$ 给出 $\kappa_r=2^{r(r+1)}L$，\texttt{round98sources} 再由这些输入导出径向条件。因此，接受笔记列出的结构与精确恢复输入即可得到 $C=L$、$r_0=4$ 的R。本次修订已检查后端解析推导F/J $\Rightarrow$ R；近期隔离结果已建立双常数粗界所需的实际恢复、尾界及极值图入口估计，但尚未建立这条通向R的路线所需的零菱形结构窗口、深度归纳结论，以及实际序列的固定赋值展开F与共同全倍数J。抽象解析推导仍须与其对实际极值序列的应用分开表述。准确说法是“笔记中存在证明路线，但本文未对整条路线完成独立认证”，而不是“笔记没有论证”。因此附录C、D的实际应用仍显式保留R；附录A不依赖此长链。

\section*{附录C. 伸缩关系的排除证明}
函数芽称为D-finite，若满足非零的多项式系数线性微分方程。本附录先说明自然边界，再使用R中的精确发散系数研究伸缩关系。

奇点方法在非D-finite性研究中已有系统应用，参见Flajolet、Gerhold与Salvy [5]。这里需要区分自然边界与Mahler方程：Bell、Coons与Rowland [6] 研究Mahler函数的有理性—超越性二分及自然边界，非有理Mahler函数本身可以具有自然边界。因此，自然边界不能单独排除Mahler关系；本附录的排除使用的是R中随单位根阶数变化的首个发散导数及其系数。

\textbf{定理C.1.} 若 $G$ 满足R，则单位圆为其自然边界，且 $G$ 不是D-finite函数。

\textbf{证明.} R使每个足够高阶的原始二进根处某个有限阶径向导数对数发散，故 $G$ 不可能在该点解析延拓。这些根在单位圆上稠密，因而不能穿过任何非空圆弧延拓。D-finite函数的局部分支只能在其微分方程最高系数的有限零点处发生有限平面奇性，故不可能有此自然边界。证毕。

下面的右端条件稍强于“圆盘内解析”：清除有理系数的公共分母后，右端在盘内选定的分支可穿过单位圆除有限点之外的每一点解析延拓。D-finite右端满足这个条件。

\textbf{定理C.2（混合伸缩关系的必要最高项）.} 设 $\mathcal M$ 为有限个不同正整数组成的集合。考虑非平凡关系

\[
\sum_{m\in\mathcal M}Q_m(z,\Theta)[G(z^m)]=F(z),
\]
\[
Q_m(z,\Theta)=\sum_{j=0}^{d_m}p_{m,j}(z)\Theta^j,
\]
其中 $p_{m,j}\in\mathbb C(z)$，$p_{m,d_m}\ne0$，且 $F$ 满足上述右端条件。系数位于算子左侧，$\Theta$ 作用于整个复合函数。写

\[
m=2^{s_m}t_m,\qquad t_m\in2\mathbb Z+1,\qquad w=\max_{m\in\mathcal M}(d_m+2s_m).
\]
对每个奇数 $t$，令 $I_t=\{m\in\mathcal M:d_m+2s_m=w,\ t_m=t\}$。若 $G$ 满足R，则

\[
\sum_{m\in I_t}2^{-s_m^2}t^{-2s_m}p_{m,d_m}(z)=0
\]
是有理函数恒等式。这只是该关系存在的必要条件，并非充分条件；最高项抵消不保证低阶项也能满足原关系。

\textbf{证明.} 清分母，使全部系数为多项式；公共分母只乘在算子左侧，不与 $\Theta$ 交换。固定充分大的 $r$，使 $r-s_m\ge r_0$ 且 $K=2r-w\ge0$。取原始 $2^{r+1}$ 次根 $\omega$。由于 $\omega^m$ 的阶为 $2^{r+1-s_m}$，$G(\omega^m y)$ 的首次发散的导数阶为 $k_m=2(r-s_m)$。

沿 $z=\omega x$，$\Theta=x\,d/dx$。对关系再求 $K$ 阶径向导数，只有满足 $d_m+2s_m=w$ 的最高微分项达到实际导数阶 $k_m$；其余项只涉及低阶导数，均有有限极限。乘积法则、Euler展开和复合求导中唯一的最高项系数在端点为 $p_{m,d_m}(\omega)m^{k_m}$。因此该项的首个对数发散项的系数为

\[
C\,2^{(r-s_m)^2-1}m^{2(r-s_m)}p_{m,d_m}(\omega)
=C\,2^{r^2-1}t_m^{2r}2^{-s_m^2}t_m^{-2s_m}p_{m,d_m}(\omega).
\]
所有乘子在端点的差为 $O(1-x)$，与对数相乘仍有界。除去右端有限个例外点，其同阶导数有限，故这些首个对数发散项的系数之和为零。

对每个固定 $r$，于是得到一个关于 $\omega$ 的多项式等式；该多项式次数由原系数次数控制，与 $r$ 无关。原始 $2^{r+1}$ 次根有 $2^r$ 个。当 $r$ 足够大时，除去有限例外仍超过这个固定次数，故此多项式恒为零。

按不同奇数部分归组，记 $T_t(z)=\sum_{m\in I_t}2^{-s_m^2}t^{-2s_m}p_{m,d_m}(z)$，则对所有充分大的整数 $r$ 有 $\sum_t t^{2r}T_t(z)=0$。选取与不同 $t$ 个数相同的连续 $r$，系数矩阵是以不同 $t^2$ 为基的Vandermonde矩阵乘一个非零对角矩阵，因而可逆。于是每个 $T_t=0$。本证明逐个固定 $r$ 取极限，随后只作有限代数消元，不交换增长深度与径向极限。证毕。

\textbf{推论C.3.} 在R下，不存在非平凡的有限关系 $\sum_m p_m(z)G(z^m)=F(z)$，其中系数有理且右端满足定理C.2的条件。

\textbf{证明.} 此时所有 $d_m=0$，最高权重组具有相同的最大 $s_m$。其中不同 $m$ 的奇数部分不同，每个非空 $I_t$ 只有一项，与非零最高系数矛盾。证毕。

\textbf{推论C.4.} 若任意两个伸缩量之比都不是2的整数幂，则即使允许任意有限阶微分，定理C.2的关系也不可能成立。特别地，若整数 $b\ge2$ 不是2的幂，则不存在这种有限阶线性微分—Mahler关系 $\sum_jQ_j(z,\Theta)[G(z^{b^j})]=F(z)$。

\textbf{证明.} 不同伸缩量具有不同奇数部分，所以最高权重的每个非空组均为单项。对于 $b=2^hu$、奇数 $u>1$，$b^j$ 的奇数部分为互异的 $u^j$。证毕。

同理，任何可能关系中，最高权重的每个非空奇数部分组至少有两项；它们的微分阶之差必须为 $2(s_2-s_1)$。全部最高微分阶的跨度不超过1时，关系不可能成立。对于两项关系 $p(z,\Theta)G(z)-q(z,\Theta)[G(z^2)]=F$，必要条件准确为 $\deg_\Theta p=\deg_\Theta q+2$ 及 $q_{\mathrm{top}}=2p_{\mathrm{top}}$。这只是必要条件。

\section*{附录D. 数字线性表示与线性递推}
序列称为P-recursive，若最终满足非零、有限阶、多项式系数线性递推。普通生成函数D-finite与系数P-recursive等价[3]；允许有限个初始项例外。

对序列 $(b_n)_{n\ge0}$，固定整数基数 $q\ge2$，考虑所有子序列 $n\mapsto b_{q^en+t}$，其中 $e\ge0$、$0\le t<q^e$。如果这些序列的复线性张成空间有限维，称该序列为 $q$-regular[4]。

\textbf{定理D.1.} 若 $G(z)=\sum_{n\ge0}b_nz^n$ 满足R，则其系数序列在任意固定整数基数 $q\ge2$ 上均非regular；有限项修改也不能改变这个结论。

\textbf{证明.} 反设核空间维数为 $D$。选基向量 $v(n)$，按模 $q$ 抽取子序列的操作保持该空间不变，因此存在常矩阵 $M_j$，满足 $v(qn+j)=M_jv(n)$，并有 $b_n=\ell v(n)$。向量生成函数 $V$ 满足

\[
V(z)=\left(\sum_{j=0}^{q-1}z^jM_j\right)V(z^q),\qquad G(z)=\ell V(z).
\]
这在形式幂级数意义成立，包含 $n=0$；不需要矩阵可逆。将 $G(z),G(z^q),\ldots,G(z^{q^D})$ 展开到同一个向量 $V(z^{q^D})$，得到 $D+1$ 个属于 $\mathbb C(z)^D$ 的行向量。线性相关给出一个非平凡、无导数的有限伸缩关系，违反推论C.3。若只改有限项，生成函数只差多项式，右端仍为有理函数，矛盾依旧成立。证毕。

规范数字表示中允许首个非零数字初始化为固定向量，只在向量生成函数中增加有限多项式。增广一个常数方向后同样消元。因此定理也排除固定维度、常矩阵更新及固定线性读出的精确数字模型；它不排除随输入变化的状态维数、非线性更新或一般算法。

\textbf{条件性推论D.2.} 若实际 $A$ 满足R（$C=L$），则 $A$ 的自然边界为单位圆，$E$ 的自然边界为半径2的圆，且 $P(n)$ 非P-recursive。附录C的伸缩关系限制在此前提下适用于 $A$。

其中 $P(n)$ 非P-recursive由 $A$ 非D-finite得到；$P(n)=n!2^{1-n}a_n$ 的因子是超几何因子，在有限阶递推中代入并清分母，P-recursiveness在这两个序列间等价。实际 $a_n,m_n,P(n)$ 的数字表示排除则由推论A.6无条件给出，无需引用R或定理D.1。后者保留为满足R的一般序列的结构性结论。

若日后能够独立确定实际余项 $m_n$，则逐层展开正文的半规模关系可得

\[
a_n=\sum_{j=0}^{\nu_2(n)}\frac{2^{j(j+2)}}{n^{2j+2}}m_{n/2^j}.
\]
它目前只是重建恒等式：把 $m_n$ 定义为未知 $a_n$ 的差，再代回这个式子，并没有得到一个独立闭式。

\section*{附录E. 有自然边界的短公式}
正文的例子对应下面的 $C=1$。这里允许任意固定的非零复数 $C$。

固定 $C\ne0$，定义

\[
s_n=-C\sum_{j=0}^{\nu_2(n)-1}2^{j(j+1)}n^{-2j-1},\qquad S(z)=\sum_{n\ge1}s_nz^n.
\]
奇数时空和为零。对偶数分离首项得到 $s_n=-C/n+2s_{n/2}/n^2$，故准确有

\[
\Theta^2S(z)-2S(z^2)=-\frac{2Cz^2}{(1-z^2)^2}.
\]
另一方面，固定 $r$，将定义分成 $j<r$、$j=r$ 和 $j>r$。后段相邻存在项之比至多 $1/4$，因此尾系数为 $O_r(n^{-2r-3})$。令 $q=2^{r+1}$。在原始 $q$ 次根处，低层项只产生非共振的多重对数，首共振项为 $-\kappa_r\operatorname{Li}_{2r+1}(z^q)/q^{2r+1}$，其中 $\kappa_r=C2^{r(r+1)}$。沿 $z=\omega x$ 求 $2r$ 阶导数，最高链式因子为 $q^{2r}$，故对数系数为 $\kappa_r/q=C2^{r^2-1}$；低阶链式项有限，尾部该阶导数绝对可和。由此得到条件R，因此 $S$ 具有自然边界。

同时 $-s_n/C$ 是一个至多 $\nu_2(n)$ 项的有理数和。奇数 $n$ 时直接输出零；当 $\nu_2(n)\ge1$ 时，公分母可取 $n^{2\nu_2(n)-1}$，分子、分母的位长均为 $O((\log n)^2)$，从而可在输入位长的多项式时间准确计算。这里准确输出的是有理数 $-s_n/C$，或可直接取 $C=1$。这个例子同时具有自然边界和短微分—二进递推。其系数趋于零，不是实际的 $a_n$；不能把它的有理右端移植到定理A.4所处理的实际余项。

正文还提到 $1/n$ 的生成函数。由几何级数逐项积分可得 $\sum_{n\ge1}z^n/n=-\log(1-z)$。绕 $z=1$ 延拓一圈，对数的值改变 $2\pi i$，反复绕行得到无限多个分支；代数函数仅有有限个分支，因此该生成函数不是代数函数。

\section*{研究状态与许可}
本文为可供审阅的研究稿，尚未经过外部同行评审。AI工具参与了推导整理、翻译、排版与内部一致性核对；作者负责审阅最终内容。

独立的隔离整合快照于2026年11月10日08:20:17 UTC通过全部557模块的全新编译及全部7,907个声明的类型与公理审计。其中包含原有413个未改动模块及144个新增模块。各项目模块按拓扑顺序直接调用Lean编译至隔离对象目录，再检查导入内核声明清单的一致性并审计全部声明的类型与公理；这不是一次新的Lake构建。未复用历史项目对象；同一轮全新核验中断后的检查点经哈希核对后续跑。独立审计未发现非标准公理或来源、覆盖错误。该结论核验的是这一精确整合快照，不表示GitHub或规范源码树已经更新发布。复制的 \path{lean/verification} 目录仍为413模块历史记录；全新整合另有独立核验记录。

除附录A.4的非代数性结论外，形式化现已覆盖实际序列 $a_n$ 与 $m_n$ 在每个整数基数 $q\ge2$ 下的非 $q$-regular 性，采用 $q$-核的复线性张成空间有限维这一定义。对 $P(n)$，形式化覆盖增长障碍以及固定有限数字集合、连续线性状态转移和固定线性输出模型的排除，包含有限维复矩阵模型；该应用尚未形式化此表示与 $q$-核定义之间的等价。

定理A.5的一般有理极限断言现已完整形式化：有理数值的 $q$-regular 序列除以下标的非负整数次幂后若收敛，则极限必为有理数。新增的 \texttt{RationalRadial}、\texttt{RationalLinearDescent}、\texttt{RationalSeriesLimit} 和 \texttt{QRegularRationalLimit} 模块提供了相应桥接论证与最终定理。\texttt{ExtremalAsymptotic} 与 \texttt{FourExponentialIndependence} 仍为显式外部输入：图论渐近定理本身以及所需的Lindemann--Weierstrass推论均未在此Lean项目中证明。$P(n)$ 的增长论证只需前一个输入。从F与共同全倍数J到径向渐近的抽象推导已形式化，但尚未对实际极值序列建立这些前提；完整条件R及附录C--E的完整后续论证尚未在此获得形式化认证。历史核验与全新整合核验记录须按上述说明区分；源文件的存在本身不能作为验证成功的凭据。

\paragraph{已核验整合中的实际粗界。}
令
\[
 L=\frac{\cosh1}{\cos1},\qquad
 U=\frac{3\pi^4+4\pi^2-32}{\pi^4+4\pi^2-32}.
\]
存在 $C\ge0$ 与 $N\ge2$，使每个 $n\ge N$ 均满足
\[
 (L-C/n)\mu_n\le P(n)\le(U+C/n)\mu_n.
\]
该上界也适用于每个 $n$ 阶竞赛图的 $H(T)$。
记 $s=S\mathbf1$、$\tau(T)=\sum_i(s_i/(n-1))^2$，则存在
$D\ge0$ 及共同规模阈值，使每个实际极值图均满足
\[
 \tau(T)\le2\log(U/L)+D/n.
\]
最终端点不再外加路径比较、逐点得分上界、对数能量窗口或谱上界假设。
这些端点已纳入上述通过核验的隔离整合快照。这里 $U>L$。
上述结论给出归一化得分能量有界，不给出 $\tau=O(1/n)$、
常数为 $L$ 的尖锐渐近定理、结构分类，亦不证明实际序列的F/J或R。

已验证的上游引理还包括有限有序消去、颜色词的最短交换与零菱形竞赛图修复、实际竞赛图矩阵的向量弱主序，以及所有非负整数次幂的Frobenius范数比较。全一向量源能量级数及实际预解式恒等式已在 $q\ge0$、$q(n-1)^2<1$ 的实参数范围内证明，包括 $n>0$ 时的 $q=1/n^2$；收敛性与双边逆均在证明中建立。精确偶子集卷积S1与零菱形两永久式平均直接采用本仓库的实际哈密顿路计数。实际强连通分量的有序块重构还给出有限小缺陷界：当 $n\ge4$、$D>0$、$24D\le n^2$ 时，$6n^2D\le\lVert T_n^3\rVert_F^2-\lVert S^3\rVert_F^2$。这里 $S$ 是原图的反对称符号矩阵，$T_n$ 是传递竞赛图的符号矩阵，$D$ 是无标号四顶点菱形数。

对实际零菱形竞赛图，形式化现已覆盖固定复圆盘。有限核恒等式为
\[
\det(I+aS)=\frac{(1+a)^n+(1-a)^n}{2}.
\]
虚参数归一化核在复紧集上一致趋于 $\cos w$；当 $|w|\le R$，且 $n>0$、$R/n\le1/2$、$R^2/n\le1$ 时，误差不超过 $2e^R R^2/n$。对每个固定的 $0\le R<\pi/2$，还有一个不依赖图或圆盘参数的阶数阈值，使核真正双边可逆，并满足
\[
\left\|(I+(w/n)^2S^2)^{-1}\right\|_{2\to2}\le C_R^2,
\qquad C_R=\frac{2(1+R)e^R}{\cos R}.
\]
此界来自显式逆矩阵条目及欧氏算子范数估计，而非仅由行列式下界推出。对实际的实得分向量 $b=S\mathbf1$，若另有能量条件 $\lVert b\rVert_2^2\le Kn$，则
\[
\frac{\left|b^{\mathsf T}(I+(w/n)^2S^2)^{-1}b\right|}{n^2}
\le\frac{KC_R^2}{n}.
\]
得分能量恒等式，以及Hermitian内积与普通双线性源项之间的转换，均已形式化。零菱形条件本身不推出上述线性能量界。

对按零点分支定义的高斯—源项因子，令
\[
 \beta_n(z)=1-z^2/n^2,\quad k_n(z)=z\exp[-\tfrac12\log\beta_n(z)],\qquad
 G_n(z)=\frac{\exp[-\tfrac n2\log\beta_n(z)]}{H_n(k_n(z))},
\]
其中对数确定零点处的分支，$H_n(w)=((1+iw/n)^n+(1-iw/n)^n)/2$。采用实际得分源项
$\Psi_S(w)=b^{\mathsf T}(I+(w/n)^2S^2)^{-1}b$，定义
\[
 \Lambda_{v,S}(z)=\frac{vz}{1+vz/n}
 -\frac{k_n(z)^2\Psi_S(k_n(z))}{n^2(1+vz/n)},\qquad
 P_{v,S}(z)=G_n(z)e^{\Lambda_{v,S}(z)}.
\]
已验证的有限行列式恒等式给出
$G_n(z)^2\det(\beta_n(z)I+(z/n)^2S^2)=1$ 及 $G_n(0)=1$，适用于所述非零分母范围。参数分支、实际源项与组合乘积在充分大阶数下具有共同解析圆盘。固定 $0\le r<R<\pi/2$、$K\ge0$ 及任一固定导数阶数 $j$，形式化给出显式常数 $C_{j,r,R,K}$，使
\[
 \sup_{|z|\le r}\left|\frac{d^j}{dz^j}
 \left(P_{v,S}(z)-\frac{e^{vz}}{\cos z}\right)\right|
 \le \frac{C_{j,r,R,K}}n.
\]
结论对充分大的 $n$ 成立，并且统一适用于满足 $\|S\mathbf1\|_2^2\le Kn$ 的实际零菱形竞赛图和所有复参数 $|v|\le1$。导数界由较大圆盘上的函数值误差及 Cauchy 估计推出，并非对逐点误差符号直接求导。

对原始双分支 $v^2=1$，有限分块桥接也已证明。令 $M_v=(S+vI)/n$、$Q_v=\left(\begin{smallmatrix}0&M_v\\M_v^{\mathsf T}&0\end{smallmatrix}\right)$、$B_v=I+zQ_v$。在所述非零分母及可逆条件下，$B_v$ 真正双边可逆，且原始 $2n$ 维源项满足
\[
 \tfrac12\mathbf1^{\mathsf T}(B_v-I)B_v^{-1}\mathbf1=\Lambda_{v,S}(z),
 \qquad G_n(z)^2\det B_v=1.
\]
左侧源项独立地从分块矩阵定义；与消元后表达式的等价由分块方程及反对称性推出。上文对复参数 $|v|\le1$ 的推广只涉及按分支定义的因子，这些分块恒等式仍要求 $v^2=1$。

第一个有限配对桥接也已形式化：对异色条目为 $M_{ij}$、同色条目为零的分块协方差，两个带标号的 $n$ 元集合上全部完美配对的权重和等于 $\operatorname{per}M$。每条无向边在权重乘积中恰取一个代表端点；含同色边的配对权重为零，剩余二部配对与排列之间建立了真正的双射，而非直接把配对和定义成永久式。空情形 $n=0$ 亦覆盖，此时配对和与永久式均为一。这是交换半环上的有限代数恒等式，不是高斯测度期望定理。

有限配对恒等式现也适用于原矩阵在实际左右顶点子集 $I,J$ 上的限制。若 $|I|\ne|J|$，配对和为零；若 $|I|=|J|=k$，则配对和等于按显式有限枚举得到的所选 $k$ 阶子矩阵的永久式。协方差先限制到这些子集，再取配对和，因此结论直接涉及原矩阵条目，并非另行假设一个无关的方阵。

全局平方自由步骤现也已证明：在多元多项式上，用实际子集配对权重显式定义线性系数收缩映射。展开 $\prod_{x\in[n]\sqcup[n]}(1+X_x)$，再以异色协方差 $zM$ 作该收缩，得到冻结多项式
\[
 F_M(z)=\sum_{k=0}^{n}F_k(M)z^k,\qquad
 F_k(M)=\sum_{|I|=|J|=k}\operatorname{per}M[I,J].
\]
缩放关系及按子集大小重新分组均已证明，并有 $F_0(M)=1$、次数不超过 $n$。这一收缩在平方自由单项式上对应已证明的有限配对权重，但其定义忽略含重复变量的单项式。因此它尚不是重复指标的Wick泛函，不能据此直接处理指数或对数级数。

另一个完整多项式Wick泛函现已保留重复标签：对指数向量为 $d$ 的单项式，配对位置取为各个独立出现 $(i,a)$，其中 $0\le a<d_i$，而非仅取 $d$ 的支集。单项式的值由这些出现位置的全部完美配对求和，再线性延拓到多元多项式。对对称协方差，任意出现位置枚举都给出相同值，包括混合两种颜色的重排。此外，出现位置配对矩的完整四阶公式已证明，包括
\[
 m_C(x,x,y,y)=C_{xx}C_{yy}+2C_{xy}^2,
\]
无需对角为零的假设，也未丢弃对角贡献。

完整多项式Wick泛函现已证明在平方自由基单项式上与先前收缩一致。桥接采用出现位置与所选顶点子集之间保标签的显式等价。因此，对真实乘积 $\prod_x(1+X_x)$ 及缩放后的分块协方差，完整泛函也给出冻结多项式 $F_M(z)$。这只是平方自由子空间上的兼容，并非两个映射在含重复变量的一般多项式上相等。

对任意有限出现位置集及对称协方差，全阶配对递归已证明：
\[
 m_C(V)=\sum_{s\in V\setminus\{o\}} C(o,s)\,
 m_{C|_{V\setminus\{o,s\}}}(V\setminus\{o,s\}).
\]
证明真正构造了配对与“指定位置的配偶加剩余配对”之间的双射，并推导边权乘积的分解。奇数个位置的矩为零，空矩为一。这些是有限代数结论；

对常协方差 $c$，递归给出真实配对矩 $m_c(2k)=(2k-1)!!c^k$、$m_c(2k+1)=0$，空矩为一。配对计数及其阶乘归一化均已证明，而非作为假设输入。显式的出现位置等价将这些矩与完整单变量多项式Wick泛函 $\mathcal W_c(X^m)$ 联系起来。在交换 $\mathbb Q$-代数上，由此得到形式幂级数恒等式
\[
 \sum_{m\ge0}\frac{\mathcal W_c(X^m)}{m!}t^m
 =\exp(ct^2/2).
\]
右侧是经缩放和次数加倍的形式指数级数，不声称解析收敛或概率测度解释。

标量恒等式已推广到任意有限线性源 $\ell_h(X)=\sum_i h_iX_i$ 与对称协方差 $C$。通过带标签出现位置展开、真实配对端点双射及每个固定配对的因式分解，证明
\[
 \mathcal W_C(\ell_h^m)=m_{h^{\mathsf T}Ch}(m),\qquad
 \sum_{m\ge0}\frac{\mathcal W_C(\ell_h^m)}{m!}t^m
 =\exp\!\left(\tfrac12 h^{\mathsf T}Ch\,t^2\right).
\]
第二式是交换 $\mathbb Q$-代数上的形式幂级数恒等式，来自实际的重复指标多项式泛函。它要求协方差对称，不要求正定或高斯概率测度。

完整配对泛函还满足代数性的分部积分恒等式
\[
 \mathcal W_C(X_iP)=\sum_j C_{ij}\mathcal W_C(\partial_jP).
\]
该式适用于有限标号集、交换半环上的任意对称协方差及任意多项式 $P$。证明从真实配对中的出现位置删除出发：单项式中的重数 $d_j$ 来自对该标号的 $d_j$ 个出现位置求和，$d_j=0$ 时相应纤维为空；再由多项式归纳得到上述恒等式。因此重复标号并未被丢弃，矩递推也不是额外假设。这里的“分部积分”仅指形式代数恒等式，并不涉及概率积分。

对特征零域上的对称矩阵 $K,A$，混合形式级数
\[
 E(u,X)=\exp(u h^{\mathsf T}X)\exp(-u^2X^{\mathsf T}AX/2)
\]
由有限系数和定义。逐系数应用实际Wick泛函，得到 $Z_W=\mathcal W_K(E)$、$G_i=\mathcal W_K(X_iE)$ 及 $H_{ij}=\mathcal W_K(X_iX_jE)$。对多项式变量求导并使用代数分部积分，可得两级 Ward 方程
\[
 B G=u(Kh)Z_W,\qquad B H=KZ_W+u(Kh)G^{\mathsf T},\qquad B=I+u^2KA.
\]
形式矩阵 $B$ 已证明可逆，且具有真实的双侧逆。令 $r=B^{-1}Kh$、$D=B^{-1}K$。求解上述实际 Ward 方程，不需要除以 $Z_W$，即可得到
\[
 G=urZ_W,\qquad H=(D+u^2rr^{\mathsf T})Z_W.
\]
关于参数的导数也已逐系数证明，并有 $Z_W(0)=1$：
\[
 Z_W'=q(u)Z_W,\qquad
 q(u)=u h^{\mathsf T}r-u\sum_{i,j}A_{ij}(D_{ij}+u^2r_ir_j).
\]
初值来自实际空配对矩。完整的 Jacobi 求导公式、逆矩阵求导及形式标量微分方程的唯一性现已推出一般混合源恒等式
\[
 \begin{aligned}
 Z_W(u)&=\Delta(u)\exp(\Phi(u)),&
 \Phi(u)&=\frac{u^2}{2}h^{\mathsf T}B(u)^{-1}Kh,\\
 \Delta(u)^2\det B(u)&=1,& \Delta(0)&=1.
 \end{aligned}
\]
该结论适用于有限标号集、具有有理数代数结构的特征零域上的任意对称矩阵 $K,A$ 及任意源向量 $h$，不要求协方差 $K$ 可逆。行列式分支由指数为 $-1/2$ 的通用二项级数代入 $\det B-1$ 构造，形式指数则代入 $\Phi$；两处代入对象的常数项均已证明为零。证明从置换行列式导出完整 Jacobi 求导公式，在不假定 $B$ 对称的前提下证明 $B^{-1}K$ 对称，并核对候选式与实际Wick级数满足同一微分方程和初值。因此该式对应原有的有限Wick系数，包含重复指标，而非预设的高斯公式。此前的秩一二次分支亦有独立证明。

实际Wick级数的 $z=u^2$ 分级下降现已证明：下降级数的 $z^n$ 系数取自原级数的 $u^{2n}$ 系数，所有奇数次系数为零，次数加倍展开重构原级数。这是系数抽取，而非在 $u=1$ 处求值。矩阵 $B_z=I+zKA$、其真实逆、行列式分支和源指数均通过展开同态输运。

取 $A=I$、$h=\mathbf1$，源项已识别为真实矩阵二次式
\[
 \frac12\mathbf1^{\mathsf T}(B_z-I)B_z^{-1}\mathbf1.
\]
所得形式行列式–指数恒等式也已特化到原始双倍竞赛图协方差 $Q_v$，其对称性已证明，适用于任意有限阶数及任意复数 $v$。这不扩大另一个要求 $v^2=1$ 的块矩阵化简结论。

行列式因子现已有局部解析桥。以迭代导数除以阶乘定义的实际Taylor系数，已证明与有限算术、行列式和真实矩阵逆相容。对任意有限复矩阵 $Q$，$I+zQ$ 及其行列式、逆矩阵在零点的Taylor级数均与相应形式对象一致；这一矩阵结论不要求 $Q$ 对称或正定。再由归一化逆平方根分支的唯一性，对无菱形竞赛图且 $v^2=1$，现有解析高斯因子的Taylor级数已识别为规范形式行列式分支。其解析性、零点取值一及局部平方恒等式均由实际定义证明，而非外加假设。

原始源与指数的Taylor恒等式现也已证明。对任意有限复矩阵 $Q$，实际逆二次源在零点附近解析、零点取值为零，且具有形式原始源的系数。Taylor求导、解析链式法则及形式微分方程唯一性给出其指数的识别，不需要交换无穷求和。结合归一行列式分支，得到真正局部收敛的展开，其系数正是此前定义的下降完整Wick系数。对 $n>0$、无菱形竞赛图且 $v^2=1$，原始源已识别为原来的块矩阵表达式，再经有限源消元，将该展开识别为现有乘积 $P=G\exp\Lambda$ 的展开。因此具有实际Wick系数级数的是原有乘积，而非另行定义的候选函数；$Q_v$ 中已有的归一化保持不变。

收敛现已推广到具有共同规模阈值的整个闭内圆盘。固定 $0\le r<R<\pi/2$ 后，存在同一个 $n$ 阈值，使任意无菱形竞赛图、任意 $v^2=1$ 和任意 $|z|\le r$ 都满足：实际Wick级数之和等于原有乘积。此结论不需要得分能量假设。这里证明的是上述统一量词下每一点的级数求和，不是尚未建立的一致收敛速率。

以 $c_k$ 记实际Wick系数，其估计现已有经证明的图统一上界。对 $0\le\rho<s<R<\pi/2$、$K_0\ge0$，假设得分能量不超过 $K_0n$。令 $M=\exp(s)/\cos R+C(s,R,K_0)$，其中 $C$ 为已证明的乘积误差常数；在同一规模阈值后，对所有这样的无菱形竞赛图及 $v^2=1$，有
\[
 |c_k|\le M/s^k,\qquad \sum_{k\ge0}|c_k|\rho^k\le\frac{M}{1-\rho/s},
\]
\[
 \left|P(T,v,z)-\sum_{k<m}c_kz^k\right|
 \le\frac{M(\rho/s)^m}{1-\rho/s}\quad(m\ge0,\ |z|\le\rho).
\]
球面上界来自实际乘积估计，而非额外假设。最后一式是解析Wick核心的Taylor余项，不是有限匹配多项式或永久式的真实尾项估计。

对任意多项式插入 $F$，插入Ward系统及其经真实逆得到的解也已证明。令 $V=B^{-1}K$、$r=Vh$，归一化倾斜泛函满足
\[
 \widetilde{\mathcal W}(X_iF)
 =\sum_jV_{ij}\widetilde{\mathcal W}(\partial_jF)
   +u r_i\widetilde{\mathcal W}(F),\qquad
 \widetilde{\mathcal W}(1)=1.
\]
归一化使用已证明可逆的 $Z_W$，保留重复指标。实际多项式插入的物理缩放与偶下降现已证明。对标号词 $w$，物理矩定义为 $u^{|w|}\widetilde{\mathcal W}(\prod_{i\in w}X_i)$，每个出现位置对应一个因子；首个出现位置的删除递推由多项式Ward恒等式推出。另行定义的终止 $z$ 递推使用
\[
 W=z(I+zKA)^{-1}K,\qquad \eta=Wh,
\]
并经 $z=u^2$ 展开重构上述实际矩。取 $A=I$ 时，真实双侧逆恒等式还给出 $W=zK(I+zK)^{-1}$。$W$ 与 $\eta$ 的常数项均为零。空矩为一，所有奇数次 $u$ 系数为零；下降级数的 $k$ 次系数在 $2k<|w|$ 时为零。词的任意排列不变性由实际交换多项式乘积输运得到，而非递推的假设。重复标号和空词均保留。这些是有限形式恒等式，不引入得分能量或图规模假设。

残余因子分解现已在 $u$ 分级的多项式系数级数层面证明。独立定义
\[
 r_i=(1+uX_i)\exp(-uX_i+u^2X_i^2/2),\qquad R=\prod_i r_i.
\]
多项式代数中形式指数的加法与逆恒等式给出 $ER=\prod_i(1+uX_i)$，其中 $E$ 正是 $A=I$、$h=\mathbf1$ 时的既有核心。直接系数计算得到 $[u^0]r_i=1$、$[u]r_i=[u^2]r_i=0$。残余因子不是以目标答案除以核心定义的。

对 $R=\sum_aR_au^a$，归一化插入卷积定义为
\[
 [u^n]\mathcal T(R)=\sum_{a+b=n}[u^b]\widetilde{\mathcal W}(R_a).
\]
由实际Wick线性性及已证明的可逆归一化，得到
\[
 \mathcal W\!\left(\prod_i(1+uX_i)\right)=Z_W\mathcal T(R).
\]
每一系数均为有限求和；这里既不使用Wick的乘法性，也不进行解析求和，重复标号及自配对均保留。

有限冻结多项式现已与上述实际残余联系起来。对真实子集配对权重归并后，已证明移位乘积的Wick值是既有冻结多项式之规范幂级数像的次数加倍展开；不等势子集由配对定理给出零贡献。以上述实际插入卷积为 $T_U$，按 $[z^k]T_Z=[u^{2k}]T_U$ 定义 $T_Z$。利用实际核心常数项为一的可逆消去，可证明 $T_U$ 的奇数系数为零，且次数加倍展开重构 $T_U$。于是，对由 $M$ 构成的原二部协方差，有形式幂级数恒等式
\[
 \operatorname{frozenPolynomial}(M)=\operatorname{core}(M)\,T_Z(M).
\]
其中核心是实际下降Wick级数，残余由有限插入系数及偶系数抽取定义，并非以目标多项式除以核心定义。残余的显式有限反对角系数公式及核心的行列式–源公式均已证明。空情形及常数项保留，不引入得分能量或渐近假设。

残余系数现已有局部有限的物理矩表达。实际局部 $m$ 次系数已证明为 $\gamma_mX_i^m$；将有限残余乘积按次数分配展开，构造每次出现各占一个位置的词，从而将实际残余系数接到已证明的 $W,\eta$ 矩。在 $z^k$ 次系数中，仅总次数 $m\le2k$ 的分配有贡献，更高次数的项为零。这是已证明的系数恒等式，不是活动展开的假设，也不是在 $u=1$ 处求值。

直接标量计算给出 $\gamma_0=1$、$\gamma_1=\gamma_2=0$、$\gamma_3=1/3$、$\gamma_4=-1/4$。因此非零分配中每个活动数值坐标的次数至少为三，活动坐标数不超过 $m/3$。数值坐标不是结构顶点，后续多个对数顶点可以共享同一标号。实际同标号一至四阶矩也已证明，其中
\[
 M_3(i)=\eta_i^3+3\eta_iW_{ii},\qquad
 M_4(i)=\eta_i^4+6\eta_i^2W_{ii}+3W_{ii}^2.
\]
自配对保留。实际标量残余现已证明等于 $\exp L$；$L$ 独立指定为零至二次系数为零、$k+3$ 次系数为 $(-1)^k/(k+3)$ 的级数，等式由实际残余的微分方程及形式唯一性推出。由 $L=u^3H$，其指数的 $n$ 次系数是对 $b\le\lfloor n/3\rfloor$ 的有限和，带有权重 $1/b!$ 及真实有序顶点次数分配。这是标量局部结构顶点展开，而非全局带标号的 $C_j$ 公式。

实际二部残余的前三个系数现已证明为 $1,0,0$。证明使用真实物理协方差系数 $[z]W_{ij}=K_{ij}$，以及总次数小于六的非零分配只有一个活动坐标这一事实。二部协方差的对角为零，相应低阶消去由此导出，并非靠丢弃自配对强行规定。因此对任意矩阵 $M$，既有冻结多项式和实际Wick核心的零、一、二次系数精确相等。这不是更高阶误差估计，活动指标 $j$ 也不等于 $z$ 的次数。

实际物理矩现已有真实部分匹配表示：在出现位置中选择配对子集及其无不动点对合，在该子集上取既有无序边权的乘积，在补集上取均值的乘积。有限子集展开将其和识别为完整Wick作用于仿射乘积，再由已证明的递推识别为原先定义的物理矩。另有真实双射将这些对象识别为允许不动点的对合；重复数值标签、内部配对和空集均保留。代入实际残余系数公式，得到这些部分匹配权重的有限和，而非另行定义一个替代对象。

对 $b$ 个次数至少为三的结构顶点，令 $D$ 为出现位置总数，$e$ 为配对子集基数的一半，$h$ 为未配对位置数。实际配对给出 $D=2e+h$。取 $j=e+h-b$ 后，已证明 $b\le2j$、$D\le6j$ 及 $e+h=j+b$。既有无序边集的基数现已证明等于 $e$。由于实际协方差和均值的常数项均为零，物理部分匹配权重可被 $z^{e+h}$ 整除，低于 $j+b$ 次的系数为零；重复数值标签及内部配对均保留。这是形式最低阶界，不是范数估计，$j$ 也不是 $z$ 的次数。

实际多项式系数残余乘积也已证明等于各局部对数之和的形式指数。其 $m$ 次系数可有限展开为 $b\le\lfloor m/3\rfloor$ 的和，带有精确因子 $1/b!$、有序顶点标签映射及总次数为 $m$ 的次数分配。标签允许重复，任一结构次数低于三的项均为零。这是原先定义的残余满足的恒等式，并非替换定义。

实际残余系数的全局有限结构归并现已证明。对有序顶点展开应用实际的加性倾斜泛函，不假设Wick或倾斜插入具有乘法性；再由真实的词与出现位置恒等式，将每项输运至物理部分匹配和。所得 $[z^k]T_Z$ 是对总次数 $m\le2k$、顶点数 $b\le\lfloor m/3\rfloor$、任意顶点标签、次数分配和实际部分匹配的有限和，保留 $1/b!$ 及对数系数。随后可将关于 $j$ 的有限和移至最外层。凡不满足 $b\le2j$、$m\le6j$ 或 $j+b\le k$ 的项，由已证明的系数最低阶与计数结论为零，而非人为截去。空结构、重复标签、同顶点配对及单点位置均保留；$j$ 不等于系数次数 $k$。

定量部分现已有若干带明确前提的工具。对绝对系数加权和可求和的级数，已证明卷积次乘性和有限乘积界；实际局部幂级数展开及较大圆周上的界，可推出较小半径处的加权系数界。将其用于既有部分匹配权重，若每个协方差的界为 $B$、每个均值的界为 $E$，则得到 $B^eE^h$；对全部顶点标签映射求和再产生 $|I|^b$，仍保留重复标签和环。可求和性及逐元素系数界是明确前提，并非形式恒等式自动推出的结论。

对维数为 $2n$、$n>0$ 的复矩阵 $Q$，设 $|Q_{ij}|\le C/n$、$|(Q\mathbf1)_i|\le K_0/n$、$|z|\le s$，且真实逆 $R(1+zQ)=I$ 的欧氏算子范数不超过 $L$。由复双线性行—逆—列估计，对 $W=(zQ)R$、$\eta=W\mathbf1$ 已证明
\[
 |W_{ij}|\le\frac{sC+2s^2C^2L}{n},\qquad
 |\eta_i|\le\frac{sK_0+2s^2CLK_0}{n}.
\]
在原始矩阵条件下，真实逆、解析性和系数桥现已由证明推出，而非作为假设。若 $\|Q\|_{\mathrm{op}}\le q$ 且 $sq<1$，收敛Neumann级数证明 $I+zQ$ 在 $|z|\le s$ 上可逆，逆算子范数不超过 $1/(1-sq)$。实际 $zQ(I+zQ)^{-1}$ 的坐标及其全一均值在该盘上解析，Taylor级数等于原先定义的形式物理协方差和均值。在上述逐元素及逐坐标分数条件下，Cauchy估计给出 $0\le\rho<s$ 处的可求和性和显式加权系数界；再用于既有部分匹配权重，闭合条件性的原始矩阵标签和估计，保留重复标签及环。分数能量界仍不能替代对 $Q\mathbf1$ 的逐坐标界。

对实斜对称矩阵，原始谱条件现已落实：若 $S$ 为实斜对称矩阵且 $|S_{ij}|\le1$，$v$ 为实数、$|v|\le1$，且 $n\ge24$，则 $M=(S+vI)/n$ 的实际二部块 $Q$ 满足 $\|Q\|_{\mathrm{op}}\le3/4$ 和 $|Q_{ij}|\le1/n$。证明使用实斜对称Lagrange恒等式、复化及真实块范数估计，不假设奇异值配对。再保留逐坐标分数条件 $|(Q\mathbf1)_i|\le K_0/n$，既有物理协方差和均值在半径 $9/8$ 处的加权系数和可求和，界分别为 $1025/(2n)$ 和 $1025K_0/(2n)$，外圆半径取 $5/4$。该定理不再另行假设谱、逆或解析性。

对显式定义的有限比特词图族，分数条件也已落实。每个正深度 $d$ 下，$2^d$ 个比特词上的平衡符号和选定枚举下的有符号传递局部矩阵，给出实际规范化行列界。经真实图重标号输运，对既有物理级数得到
\[
 \|W_{ij}\|_{9/8}\le\frac{1025}{2n},\qquad
 \|\eta_i\|_{9/8}\le\frac{1025(2^d+1)}{2n},
\]
其中范数表示绝对系数的加权和，可求和性已证明。前提为正深度、所选枚举及重标号、元素绝对值不超过一的实斜对称粗层矩阵、实数 $|v|\le1$ 及 $n\ge24$；重标号强制实际基数 $n=|C|2^d$。不再保留raw分数、谱、逆或解析性假设。这是对已定义的所选枚举图族的端到端结论；下文给出与规范有限指标popcount约定的精确识别。

有限结构活动 $C_j$ 现已直接由既有物理部分匹配权重定义，总次数不超过 $6j$，保留原对数系数及 $1/b!$。其系数恢复原先的残余：
\[
 [z^k]T_Z=\sum_{j=0}^{k}[z^k]C_j,\qquad C_0=1.
\]
实际配对纤维的基数为 $\binom{D}{2e}\operatorname{pairFactor}(e)$，并由 $2^D D^e$ 控制。跨所有长度的可容许有序次数分配可单射到 $D$ 的组合分拆，无需额外的长度因子。这些计数、$e+h=j+b$ 及保留的阶乘权重给出定量定理：若指标集大小为 $2n$、$n>0$、$M\ge1$，实际协方差和均值在半径 $\rho\ge0$ 处的加权系数和可求和且不超过 $M/n$，则对每个 $j\ge1$，
\[
 \|C_j\|_\rho\le\left(\frac{A(M)j}{n}\right)^j,
 \qquad A(M)=2^{15}\,6^3\exp(1)M^3.
\]
活动本身在该加权范数下可求和也已证明。估计使用 $j^b/b!\le\exp(j)$；保留 $1/b!$ 对得到 $j^j$ 次幂至关重要。实际权重及计数始终保留重复标签和环。

对上述所选枚举图族的正深度 $d$，取 $M=(1025/2)(2^d+1)$，即可在半径 $9/8$ 应用此定理。因而在已述粗层斜对称/元素界、重标号、实数 $|v|\le1$ 及 $n\ge24$ 条件下，对所有 $j\ge1$ 得到实际图活动界，不再外加分数、谱、逆或解析性假设。

对实际比值 $\rho(n,k)=n^k/(n)_k$，当 $n>0$、$k\le n$ 时，截至 $p$ 阶后的几何乘积余项非负且不超过 $\rho(n,k)(k^2/n)^{p+1}$。已完成构建给出单独活动的备用半径恢复界。裸活动 $C_0=1$ 的零误差不能用于乘上核心后的项。

正的径向--相位阶乘矩、永久式极化及AM--GM给出谱不等式
\[
 |\operatorname{per}A|\le\frac{n!}{n^n}h_n(s_1,\ldots,s_n)\qquad(n>0),
\]
适用于任意复矩阵 $A$，包括奇异矩阵；$s_i$ 是其奇异值。正辅助测度与不定形式协方差分开。对选定枚举的图族，由此推出半径 $5/4$ 的真实圆周界及
\[
 D_n^{-1}\sum_{k>\lfloor\delta n\rfloor}
 |[z^k]\operatorname{per}(J/n+zM)|
 \le B_{d,\delta,p}n^{-p},\qquad D_n=n!/n^n.
\]
这里 $\delta>0$、$p\in\mathbb N$、$n\ge24$，固定深度 $d>0$，粗层矩阵实斜对称且元素绝对值不超过一，采用指定枚举及重标号，并有实数 $|v|\le1$。常数显式且对该图类统一：圆周输入 $K=2^d+1$、元素界常数一及算子界 $3/4$ 均由图假设推出。此已核验尾界为实际永久式的无权绝对系数尾，不是高斯核心尾。

已核验的精确系数桥对复 $n\times n$ 矩阵 $M$ 给出
\[
 \begin{aligned}
 f_M(z)&=\frac{n^n}{n!}\operatorname{per}(J/n+zM)
       =\frac{\operatorname{per}(J+nzM)}{n!},\\
 [z^k]f_M&=\rho(n,k)F_k(M)\qquad(0\le k\le n,\ n>0).
 \end{aligned}
\]
$F_k$ 保留此前的子集配对定义。证明精确计算所选子式的 $(n-k)!$ 个补全，并核对子集行列次序；$k>n$ 时实际系数与 $F_k$ 分别为零，空矩阵规范化多项式为一。记 $G_M$ 为实际核心，则有限投影在 $k\le m$ 时将 $F_k(M)$ 识别为 $[z^k]\sum_{j=0}^{m}G_MC_j$，$j>k$ 的项为零。同一规范化多项式也已与无权图尾界的对象精确连接，不使用无限活动和。

此端点将固定正构造深度写为 $d+1$，$d\in\mathbb N$，保留上述图假设。原有实际核心满足
\[
 \|G_M\|_{9/8}\le\exp\!\left(\frac{1025}{2}(2^{d+1}+2)\right),
\]
且加权系数范数有限，对粗层矩阵、枚举、$n$ 及实数 $|v|\le1$ 统一。不假设无菱形，也不外加核心范数界。$G_MC_0=G_M$ 的乘子误差仍被保留。

对固定阶恢复，令
\[
 \begin{aligned}
 D_d&=2^{15}6^3e\left(\frac{1025}{2}(2^{d+1}+1)\right)^3,\\
 \alpha_d&=\min\!\left\{\frac12\log\frac{17}{16},\frac1{4eD_d}\right\},
 \qquad m=\lfloor\alpha_dn\rfloor.
 \end{aligned}
\]
对固定自然数 $p,\ell$，保留 $p\le m$，以
\[
 U_k=[z^k]U=\sum_{j=0}^{p}\sum_{a=0}^{p-j}
 \frac{H(k,a)}{n^a}[z^k](G_MC_j),\qquad k\ge0
\]
定义未截断形式级数 $U$。这里 $H(k,a)$ 是实际阶乘乘子的几何乘积系数；只去除了系数阶数截断，$j,a$ 两个求和仍有限。此前已核验的有限近似式为 $A(z)=\sum_{k=0}^{m}U_kz^k$，具有 $|z|\le1$ 上的实际永久式值误差界。

记 $P_k=[z^k]f_M$。加权端点证明可求和性及
\[
 \sum_{k=0}^{\infty}k^\ell\lvert P_k-U_k\rvert
 \le E_{d,p,\ell}n^{-p-1}.
\]
求和覆盖全部 $k$，而非仅有限窗口；$\ell=0$ 时 $k=0$ 的权重也为一。每个固定矩 $\sum_k k^\ell|U_k|$ 还单独证明有限。常数显式且只依赖固定 $d,p,\ell$，对所述图类的 $n$、粗层矩阵、枚举及 $v$ 统一。除 $n\ge24$ 和 $p\le m$ 外，没有新增依赖 $\ell$ 的阶数门槛；不声称对增长的深度或阶数统一。

证明保留零活动核心的乘子误差与真实 $\rho(n,k)$ 加权高活动和。在 $m$ 以上，由实际永久式的次数界将其更高阶无权尾界转为加权尾界，另用矩估计控制未截断近似式的尾部。两类尾误差均计入 $E_{d,p,\ell}$，并实际证明可求和性，不把发散和当作范数。

令下降阶乘 $(k)_\ell=k(k-1)\cdots(k-\ell+1)$，其中 $(k)_0=1$。系数接口还证明
\[
 \left|\sum_{k=0}^{\infty}(k)_\ell P_k
       -\sum_{k=0}^{\infty}(k)_\ell U_k\right|
 \le E_{d,p,\ell}n^{-p-1},
\]
并证明近似式一侧各项可求和。

解析导数桥也已核验。它对同一个未截断级数的求值函数
\[
 \mathcal U(z)=\sum_{k=0}^{\infty}U_kz^k
\]
建立导数识别。其 $17/16$ 半径加权系数范数有限，由 $G_MC_j$ 已有的 $9/8$ 范数、$0\le H(k,a)\le k^{2a}$ 及比值 $17/18$ 的备用半径矩界推出。这里只对 $j,a$ 作有限求和，不假设无限活动收敛，也不以单位半径矩代替正的半径间隙。

在 $1$ 的开邻域内逐项微分，对每个固定自然数 $\ell$ 得到
\[
 \mathcal U^{(\ell)}(1)=\sum_{k=0}^{\infty}(k)_\ell U_k,
 \qquad f_M^{(\ell)}(1)=\sum_{k=0}^{\infty}(k)_\ell P_k.
\]
多项式作为级数的求值已识别为普通多项式求值；规范化多项式先全局识别为实际函数 $\operatorname{per}(J+nzM)/n!$，再取导数。两侧解析识别均不再作为调用方前提。外半径及导数身份不要求阶数窗口条件；误差定理仍保留 $p\le m$，并给出
\[
 \left|f_M^{(\ell)}(1)-\mathcal U^{(\ell)}(1)\right|
 \le E_{d,p,\ell}n^{-p-1}.
\]
常数正是此前加权误差的常数。实际图假设及 $n\ge24$ 不变，没有新增依赖 $\ell$ 的阶数门槛。这是 $z=1$ 处的解析 \texttt{iteratedDeriv} 身份，而非对误差符号作形式微分。

规范数值编号现已精确识别。递增二进制编号以最新一位为最低位；在 $0\le i<2^r$ 上，原递归符号恰为 $(-1)^{\operatorname{popcount}(i)}$。整个张量及其不同粗层块之间的项均与已有反复倍增族一致。内部枚举反转以匹配原有序核，物理编号仍递增；深度一块保留 $J_0=\left(\begin{smallmatrix}0&-1\\1&0\end{smallmatrix}\right)$。在连续编号 $i_{\rm fine}+2^r i_{\rm coarse}$ 下得到精确矩阵等式，不外加重标号或切换等价前提。

因此，同一恢复结论适用于规范数值图族：粗层矩阵 $B$ 实斜对称、$|B_{ij}|\le1$，深度固定为 $d+1$，实数 $|v|\le1$，且 $n=q2^{d+1}\ge24$，其中 $q$ 为粗层阶数。半径与导数身份对每个 $p$ 成立；误差接口仍要求 $p\le\lfloor\alpha_dn\rfloor$，保留同一 $E_{d,p,\ell}$，没有新增依赖 $\ell$ 的阶数门槛，也不引入核、分数、逆或解析识别前提。这不证明极值分类，也不证明任意 $B$ 所得图无菱形。

另一有限投影计算取复 $n\times n$ 矩阵 $S$ 及独立复参数 $f,b$。记 $A=fS+bI$、$B=-fS+bI$、$D=(1-b^2)I+f^2S^2$、$K=D^{-1}$。若 $D$ 可逆，则实际双块系统满足
\[
 \mathcal B=\begin{pmatrix}I&A\\B&I\end{pmatrix},\qquad
 \mathcal B^{-1}=\begin{pmatrix}K&-AK\\-BK&K\end{pmatrix}.
\]
对 $n>0$，把两个非对角块中的 $bI$ 换成 $b\Pi$，$\Pi=I-J/n$，得到显式两色分解 $\mathcal B_\Pi=\mathcal B+UV$。令 $G=I_2+V\mathcal B^{-1}U$、$T=\mathcal B^{-1}UG^{-1}V\mathcal B^{-1}$。若 $G$ 也可逆，则实际逆与协方差核满足
\[
 \mathcal B_\Pi^{-1}=\mathcal B^{-1}-T,\qquad
 W_\Pi-W=T,\qquad \eta_\Pi-\eta=T\mathbf1.
\]
原有标量源之差恰为 $\mathbf1^{\mathsf T}T\mathbf1/2$。在 $f=a$、$b=av$ 时，这些核与旧实际协方差精确识别。不除以 $b$；$b=0$ 时有 $G=I_2$、$T=0$。通用恒等式仍保留 $D,G$ 可逆前提；在固定圆盘内，这些前提现已导出：若 $n\ge24$，$S$ 实斜对称且 $|S_{ij}|\le1$，$v$ 为实数且 $|v|\le1$，则对 $|z|\le5/4$ 取 $f=z/n$、$b=zv/n$，实际原始及投影系统均可逆，两者欧氏算子逆范数均至多16，并由此推出 $D,G$ 可逆。对实数 $|\sigma|\le1$ 及复数 $|t-1|\le1/8$，代换 $z=tn/(n+\sigma)$、$f=t/(n+\sigma)$、$b=f\sigma$ 落在该圆盘内，给出相同的投影逆界。不声称 $D^{-1}$ 或 $G^{-1}$ 的单独数值界。未乘 $z$ 的投影交互算子满足 $Q_\Pi\mathbf1=(s/n,-s/n)$、$s=S\mathbf1$；投影消去常参数项，并未消去实际比分源。这个粗逆界不要求分数前提，也不提供更精细的逐元投影或活动改核误差。

固定深度图族在 $1$ 处的导数识别已完成，但不声称新的更大圆盘统一恢复误差或任意矩阵、全部允许半径的raw恢复定理。无限活动收敛、增长深度/阶数统一、编辑Lipschitz余项及Euler算子闭式多项式身份仍不在该结果内。量化投影恢复及结构分类仍不在该结果内。另行完成的隔离研究现已建立实际极值图入口及上述粗界。

另两个模块给出全局入口论证G1/G2的有限输入。对任意实际竞赛图，记符号矩阵为 $S$、行分数为 $s_i$、能量为 $E=\sum_i s_i^2$；删除顶点集 $D$，$k=|D|$，令 $E_{-D}$ 为 $D^c$ 上实际主子矩阵的行和平方能量。已核验
\[
 \begin{aligned}
 E-kn^2-2k\sqrt{nE}&\le E_{-D},\\
 E_{-D}&\le2E+2(n-k)k^2.
 \end{aligned}
\]
删集大小任意，允许 $k$ 随 $n$ 增长。若 $L\ge0$ 且 $E\le Ln$，则 $E_{-D}\le2(L+1)n(k+1)^2$，不要求逐坐标分数有统一常数界。

第二个有限恒等式针对预缩放。设实斜对称矩阵 $S$ 满足 $n\ne1$、$\sum_jS_{ij}=(n-1)a_i$ 及 $|a_i|<1$，记 $v_i=a_i/(1-a_i^2)$、$w_i=a_i^2/(1-a_i^2)$、$V=\sum_iv_i$、$W=\sum_iw_i$，向量记为 $\mathbf v=(v_i)_i$、$\mathbf w=(w_i)_i$。实际规范化总质量超额满足
\[
 \begin{aligned}
 \mathfrak M&=\frac1{n-1}\sum_{i,j}
 \frac{1-\delta_{ij}+S_{ij}}{(1+a_i)(1-a_j)}-n,\\
 \mathfrak M&=\frac{W^2-V^2+W+2\mathbf w^{\mathsf T}S\mathbf v}{n-1}.
 \end{aligned}
\]
对角项 $+W$ 及有序交叉项 $\mathbf w^{\mathsf T}S\mathbf v$ 均保留。实际竞赛图特化取 $a_i=s_i/(n-1)$、$n\ge2$。对 $n\ge2$、满足 $|S_{ij}|\le1$ 及相同行分数方程的实斜对称矩阵，量化扩展要求 $|a_i|\le a_0$、$0\le a_0<1$。记 $\tau=\sum_i a_i^2$，已证明
\[
 |\mathfrak M|\le\frac{4}{(1-a_0^2)^2}
 \left(\frac{\tau+\tau^2}{n}+\frac{\tau\sqrt\tau}{\sqrt n}\right).
\]
这里 $\mathfrak M$ 为总质量减去 $n$，不是相对误差；实际图满足 $\tau=E/(n-1)^2$。不外加线性能量界或 $\tau$ 的增长率前提；证明导出 $\sum_i a_i=0$ 的抵消，并估计实际有序交叉项。

另设实 $n\times n$ 矩阵 $Q$ 的总质量恰为 $n$。任意选取行集 $I$、列集 $J$，满足 $|I|=|J|=t<n$，两集合不必相同。记 $r_i=\sum_jQ_{ij}-1$、$c_j=\sum_iQ_{ij}-1$，令 $\kappa$ 为 $\frac n{n-t}Q[I^c,J^c]$ 的总质量减去 $n-t$。则
\[
 \kappa=\frac n{n-t}\left(\sum_{i\in I,j\in J}Q_{ij}
       -\sum_{i\in I}r_i-\sum_{j\in J}c_j-\frac{t^2}{n}\right).
\]
若 $n>0$、$2t\le n$、$C,\delta\ge0$、$|Q_{ij}|\le C/n$，且各行、列残差的绝对值至多 $\delta$，则
\[
 |\kappa|\le4t\delta+\frac{2(C+1)t^2}{n}.
\]
交叉块与 $-t^2/n$ 均保留；对 $Q_{ij}=1/n$ 两者精确抵消。不假设概率独立或双随机性。实际归一化桥接现已证明。对 $n\ge2$ 且满足相同归一分数界的实际竞赛图，记 $L_0=(1-a_0)^{-1}$，以 $B_\nu$ 表示上式的 $|\mathfrak M|$ 界。实际非负矩阵 $C=(J-I+S)/(n-1)$ 总质量为 $n$，其预缩放 $C'_{ij}=C_{ij}/((1+a_i)(1-a_j))$ 满足 $C'_{ij}\ge C_{ij}/4$。故 $\sum C'=n+\mathfrak M\ge n/4$，并有
\[
 f_0=\frac n{n+\mathfrak M},\qquad \widehat C=f_0C',\qquad
 0<f_0\le4,\qquad \sum\widehat C=n,
\]
无需预设质量误差小。具体有 $|f_0-1|\le4B_\nu/n$、$|\widehat C_{ij}|\le16L_0^2/n$，两侧边际残差均至多
\[
 \begin{gathered}
 \widehat\delta=\frac{16L_0^2\sqrt\tau}{\sqrt n}+\frac{4B_\nu}{n},\\
 \|\widehat C-C\|_F^2\le
 \frac{128L_0^4\tau}{n}+\frac{512L_0^4B_\nu^2}{n^2}.
 \end{gathered}
\]
因此，对分别选取的 $t$ 元行集 $I$、列集 $J$，若 $2t\le n$，实际归一化删集质量满足 $|\kappa|\le4t\widehat\delta+2(16L_0^2+1)t^2/n$。这些是从实际矩阵导出的估计；精确总质量不等于各行列精确归一。下述局部缩放定理要求真实中心化谱间隙。后续隔离结果已建立实际粗界所需的非主子矩阵恢复及异常集估计，包括实际极值图最终满足固定逐点得分上界与归一化能量有界。它们并未建立更强的原始得分线性能量窗口或结构分类。

\subsection*{保留固定谱间隙的局部缩放}
设 $M$ 是 $m\times m$ 非负实矩阵，$m\ge1$。记
\[
 r_i=\sum_jM_{ij}-1,\quad c_j=\sum_iM_{ij}-1,\quad
 E_f=\sum_i r_i^2+\sum_j c_j^2,\quad P=J/m,\quad \Pi=I-P.
\]
假设 $M_{ij}\le C/m$、$|r_i|,|c_j|\le\varepsilon$，且
$\|\Pi M\Pi\|_2\le q$，其中 $C\ge0$、$0\le q<1$。下文算子范数均为欧氏算子范数。取
\[
 \begin{gathered}
 \beta=\tfrac32+2(C+4)+\frac{4(C+4)^2}{1-q},\qquad K=2\beta,\\
 \varepsilon_0=\min\{1,(8K)^{-1},(64K^2)^{-1},(1-q)/8\}.
 \end{gathered}
\]
若 $0\le\varepsilon\le\varepsilon_0$，则存在实向量 $x,y$，满足
$\sum_i x_i=\sum_j y_j$、$|x_i|,|y_j|\le a:=2K\varepsilon$，使
$N_{ij}=M_{ij}\exp(x_i+y_j)$ 为双随机矩阵。令 $g=1-q-4\varepsilon>0$、$S=1+\varepsilon$，同一组向量还满足
\[
 \begin{gathered}
 \sum_i x_i^2+\sum_j y_j^2\le \frac{e^{4a}}{g^2}E_f,
 \qquad N_{ij}\le e^{2a}C/m,\\
 \|N-M\|_F^2\le\frac{2CS e^{8a}}{m g^2}E_f,\\
 \|N-M\|_2\le S(e^{2a}-1),\qquad
 \|N-P\|_2\le q+S(e^{2a}-1)\le(1+q)/2<1.
 \end{gathered}
\]
这里 $E_f$ 是全部边际残差的平方和。逐坐标残差界仅给出 $E_f\le2m\varepsilon^2$，所以平方Frobenius估计一般是 $O(\varepsilon^2)$，不是 $O(\varepsilon^2/m)$；后者还需 $E_f=O(\varepsilon^2)$ 的总能量输入。

\emph{证明。} 先固定规范 $\sum x_i=\sum y_j$。实际边际映射 $F(x,y)$ 的分量是 $N$ 的各行列和减一，$F(0)=(r,c)$，零点导数为
\[
 H=\begin{pmatrix}\operatorname{diag}(\operatorname{row}M)&M\\
 M^{\mathsf T}&\operatorname{diag}(\operatorname{col}M)\end{pmatrix}.
\]
中心化把它分为具有间隙 $1-q$ 的模型及绝对行和至多 $4\varepsilon$ 的对称扰动。加入带符号规范投影，消除 $(x,y)\mapsto(x+t,y-t)$ 的一维不唯一性。密度与谱假设由此给出具有所需统一行界的实逆 $L$，并未预设缩放根已经存在。

令 $R(z)=F(z)-F(0)-Hz$。标量指数余项估计给出半径 $a$ 的无穷范数球上至多 $2S(e^{2a}-1)$ 的Lipschitz常数。上述阈值保证 $T(z)=-L(F(0)+R(z))$ 保持该闭球，且压缩常数至多 $1/2$。所得固定点同时满足真实边际方程 $F(z)=0$ 与规范条件。

欧氏估计直接作用于同一个根，不把无穷范数换算成欧氏范数。将各边际方程乘对应坐标再求和，得到
\[
 -\langle F(0),z\rangle
 =\sum_{i,j}M_{ij}(e^{x_i+y_j}-1)(x_i+y_j).
\]
对 $|s|\le2a$，有 $(e^s-1)s\ge e^{-2a}s^2$。平衡向量上模型二次型至少为 $(1-q)\|z\|_2^2$；由对称性及 $2|z_i z_j|\le z_i^2+z_j^2$，扰动能量至多 $4\varepsilon\|z\|_2^2$，不引入维数因子。故上式右端至少为 $e^{-2a}g\|z\|_2^2$。用Cauchy--Schwarz并单独处理 $z=0$，即得 $\|z\|_2\le e^{2a}\|F(0)\|_2/g$。

最后，由 $M_{ij}^2\le(C/m)M_{ij}$、$|e^s-1|\le e^{2a}|s|$ 和 $(x_i+y_j)^2\le2(x_i^2+y_j^2)$，有限求和得到
$\|N-M\|_F^2\le2CS e^{4a}\|z\|_2^2/m$，再代入根能量界。有限Schur估计把绝对行列和转成算子界。双随机性给出 $N-P=\Pi N\Pi$。阈值保证 $KS(e^{2a}-1)\le1/4$，而 $K(1-q)\ge128$，故谱误差至多 $(1-q)/512$，从而得到所述固定间隙。

该局部矩阵定理仍保留所述密度、边际误差及中心化谱假设。另行建立的实际粗界路线已消去其最终界所需的恢复与入口前提；结构分类及实际序列的F与共同全倍数J仍未证明。

另五个抽象分析模块给出F/J的若干组件。其中F、全倍数J及周期跳跃谓词只是接口，不是实际极值序列的定理。对 $b>0$、$b$ 周期序列及原始 $2b$ 次根，有限半周期配对证明首共振 $-\kappa/(2b)$。有界复系数在 $0<x<1$ 上给出每阶真实实径向导数的逐项身份；下降阶乘矩可和时，导数在1处有有限左极限。正整数项全局界 $|c_n|\le Cn^{-k-2}$ 通过 $C/n^2$ 主控，为 $h\le k$ 提供这些可和矩，不对未知函数层余项求导。

对反周期系数 $v_{n+b}=-v_n$、$b>0$，精确身份
\[
 \sum_{n\ge0}v_nz^n=\frac{\sum_{n=0}^{b-1}v_nz^n}{1+z^b},\qquad |z|<1,
\]
通过1处的有理延拓，给出未加权级数所有真实径向导数的有限左极限。独立的Dirichlet/Abel论证给出 $\sum_{n\ge0}v_{n+1}x^n/(n+1)^d$、$d\ge1$ 的有限径向极限；$d=1$ 使用自然顺序收敛，不声称绝对收敛。对每个整数 $q>0$，稀疏调和身份及实径向正则部分为
\[
 \sum_{\substack{n\ge1\\q\mid n}}\frac{z^n}{n}=-\frac{\log(1-z^q)}q,\qquad
 \sum_{\substack{n\ge1\\q\mid n}}\frac{x^n}{n}+\frac{\log(1-x)}q
 \longrightarrow-\frac{\log q}q.
\]
第一式要求 $|z|<1$，极限取 $x\uparrow1$。从F与共同全倍数J到真实径向渐近的抽象桥现已证明。对 $r\ge1$、原始 $2^{r+1}$ 次根 $\omega$，若满足固定赋值展开F及系数为 $2^{r(r+1)}L$ 的统一跳跃假设J，则每阶 $h<2r$ 导数均有有限径向极限，而 $2r$ 阶导数减去 $L2^{r^2-1}\log(1-x)$ 后在 $x\uparrow1$ 有有限极限。J作用于所有充分大的正 $2^r$ 倍数，不限于单一精确赋值类。证明包含下降阶乘修正、稀疏对数导数、Laurent装配、周期跳跃及有限赋值粘合。该端点允许 $L=0$；非零对数奇性仍需独立输入 $L\ne0$。实际极值序列的F与共同全倍数J仍未证明，完整条件R仍开放。

历史稿保留于 \path{historical/formula-restrictions}。固定版本研究笔记[7]提供了通向R的路线；命题B.1之后的依赖说明区分了已检查的解析推导与尚未全部复核的上游论证。


\paragraph{局部扩展的形式化核验。} 局部缩放端点、抽象径向桥、实际粗MainBound及归一化能量有界端点均纳入上述已核验的557模块隔离整合快照。PDF验收与发布独立于Lean核验；此处不声称已作远程发布。

© 2026 刘星辰。文本与源文件采用 Creative Commons Attribution 4.0 International（CC BY 4.0）许可。

\clearpage
\section*{参考文献}
[1] Alexander Divoux and Shouda Wang. \emph{The maximum number of Hamiltonian cycles in tournaments}. arXiv:2610.23755v1 [math.CO], 2026，2026年10月16日提交。本文引用Theorem 1.4。DOI: 10.48550/arXiv.2610.23755。固定版本全文。

[2] Sever Angel Popescu. \emph{A simple and self-contained proof for the Lindemann-Weierstrass theorem}. arXiv:2306.14352v2 [math.NT], 2023，2023年9月17日修订。本文引用Theorem 3.2。DOI: 10.48550/arXiv.2306.14352。固定版本全文。

[3] Richard P. Stanley. \emph{Differentiably Finite Power Series}. European Journal of Combinatorics \textbf{1}(2) (1980), 175–188. DOI: 10.1016/S0195-6698(80)80051-5。作者提供的正式论文。

[4] Jean-Paul Allouche and Jeffrey Shallit. \emph{The ring of k-regular sequences}. Theoretical Computer Science \textbf{98}(2) (1992), 163–197. DOI: 10.1016/0304-3975(92)90001-V。作者提供的预印本。

[5] Philippe Flajolet, Stefan Gerhold and Bruno Salvy. \emph{On the non-holonomic character of logarithms, powers, and the n-th prime function}. Electronic Journal of Combinatorics 11(2) (2005), A2. arXiv:math/0501379.

[6] Jason P. Bell, Michael Coons and Eric Rowland. \emph{The Rational-Transcendental Dichotomy of Mahler Functions}. Journal of Integer Sequences 16 (2013), Article 13.2.10.

[7] Xingchen Liu. \emph{Tournament paths research notes}. 固定提交 \path{c049e04062ed122c71f519857d54366d07b7019b}。相关笔记：\path{round96core_common_asymptotic_jets.md}、\path{round98sources_radial_log_singularities.md}、\path{round120sources_snapshot_review.md}。\url{https://github.com/LStar404/tournament-paths-research-notes/tree/c049e04062ed122c71f519857d54366d07b7019b}。

\end{document}
